\section{Finite spaces from their least neighbourhoods}
\label{sec:finite-spaces}

Consider the two-point space with opens $\varnothing,\{1\},\{0,1\}$.
Every open neighbourhood of $0$ contains $1$, whereas $1$ has a
neighbourhood missing $0$. Keeping track of this directional information
recovers the entire topology of every finite space. It also gives us small
witnesses for questions that the open-set description presents as
quantification over families.

\subsection{From a table of opens to a preorder and back}

Write $X=[n]=\{0,\ldots,n-1\}$. A subset is a bit mask $A<2^n$, a
family of masks is a mask $T<2^{2^n}$, and the complement $A^c$ is relative
to $X$. Test $T$ by checking membership of $\varnothing,X$ and checking
binary union and intersection of every pair of members. A scan starting at
$\varnothing$ and adding in each chosen member proves closure under the
union of every coded subfamily. A scan starting at $X$ proves the
corresponding intersection statement. During each scan the accumulator is
in $T$ and agrees, bit by bit, with the operation on the prefix. These are
bounded induction predicates. This also proves that complements of opens
contain $\varnothing,X$, are closed under coded intersections and finite
unions, and that taking complements again recovers $T$. Each subset used below is supplied as a mask or constructed from an
explicitly bounded decidable test.

\begin{theorem}[Finite specialization correspondence]\label{idx-finite26}
The tested topology tables on $[n]$ are in primitive-recursive bijection
with the reflexive transitive relations on $[n]$. In the convention used here
the opens are the upper sets. Refining the topology shrinks the relation.
\end{theorem}
\begin{proof}
Intersect the members of $T$ containing $x$ to obtain the open mask $N_x$.
It contains $x$, and a set is a neighbourhood of $x$ (contains an open
through $x$) exactly when it contains $N_x$. Put $x\preceq y$ when
$y\in N_x$. Reflexivity follows from $x\in N_x$. If $x\preceq y$ then
$N_x$ is an open containing $y$, so $N_y\subseteq N_x$. This proves
transitivity. For each open $U$, membership of $x$ forces $N_x\subseteq U$.
Conversely, a mask with that property is the union of $N_x$ over $x\in U$,
since each such $x$ belongs to $N_x$. Hence exactly the upper masks are open.
Starting from a preorder, the upper masks pass all the topology tests:
membership propagates along a relation edge in each union and intersection.
The principal upper set $\uparrow x=\{y:x\preceq y\}$ is upper by
transitivity and is contained in every upper set through $x$. It is
therefore $N_x$, proving the other round trip. For two tables $T,S$ on
the same carrier, $T\subseteq S$ implies $N^S_x\subseteq N^T_x$.
Conversely any mask upper for the larger relation stays upper when relation
edges are removed. This gives the refinement equivalence.
\end{proof}

The $N_x$ give a basis of at most $n$ opens. Any basis must contain each
distinct $N_x$: apply its covering property to $x\in N_x$, then use
minimality. More generally, let $\mathcal S$ be a finite table of proposed
subbasic masks. Define $x\preceq_{\mathcal S}y$ when every listed mask
containing $x$ contains $y$. This is a preorder. Its principal upper set at
$x$ is the intersection of the listed masks containing $x$, with empty
intersection $X$. Every topology containing $\mathcal S$ therefore contains
these principal sets and all their unions. Each listed mask is already
upper. Thus the preorder gives exactly the topology generated by the input
subbasis. In particular the two-point example has specialization $0\preceq1$.
On a labelled pair there are precisely four topologies: the empty and full
masks are forced, and either choice of membership for each singleton
passes the two tests.

\subsection{Interior, closure, boundary and dense sets}

\begin{theorem}[Local formulas and the closure axioms]\label{idx-finite66}
For any mask $A$ one has
\[
 x\in\operatorname{int} A\ \Longleftrightarrow\ N_x\subseteq A,
 \qquad
 x\in\operatorname{cl} A\ \Longleftrightarrow\
    (\exists y<n)\,[y\in A\ \land\ x\preceq y].
\]
Interior is the greatest open submask of $A$. Closure is the least closed
supermask of $A$, preserves binary unions and the empty set, is extensive,
monotone and idempotent. Interior preserves binary intersections and $X$,
is contractive, monotone and idempotent. The boundary consists precisely of
the $x$ for which $N_x$ meets both $A$ and $A^c$.
\end{theorem}
\begin{proof}
If $N_x\subseteq A$ and $x\preceq y$, then $N_y\subseteq N_x\subseteq A$.
Thus the first formula defines an upper mask, contained in $A$ by
reflexivity. Any open submask of $A$ that contains $x$ contains $N_x$,
proving maximality. Taking complements yields the second formula:
the finite scan for a member of $N_x$ in $A$ either produces a witness or
shows $N_x\subseteq A^c$. Closed masks are precisely lower masks.
The displayed witness formula for closure is lower by transitivity,
contains $A$ by reflexivity and is contained in every closed supermask
of $A$. A witness in $A\cup B$ supplies a witness in at
least one summand by testing its membership bit. No witness exists for
the empty mask. Inclusion preserves witnesses, and two successive
witnesses compose by transitivity. These observations give the four usual
Kuratowski axioms and monotonicity. Complement duality gives the interior
identities. Finally the boundary is $\operatorname{cl}A\cap
\operatorname{cl}(A^c)$; the witness formula proves the claim and also shows
that the boundary is closed. In particular it is empty exactly when $A$ is
both open and closed: at a point of $A$, the point itself witnesses the
first intersection test, so emptiness forces $N_x\subseteq A$, and the
same argument applies to $A^c$.
\end{proof}

Conversely, these closure axioms suffice to construct the entire topology.
Suppose a table $c$ on masks satisfies $c(\varnothing)=\varnothing$,
$A\subseteq c(A)$, $c(c(A))=c(A)$, and $c(A\cup B)=c(A)\cup c(B)$.
If $A\subseteq B$ then $B=A\cup B$ implies $c(A)\subseteq c(B)$. The
fixed masks contain $X$ and $\varnothing$ and are closed under binary
union. For any intersection $D$ of fixed masks, monotonicity puts $c(D)$
inside every factor, hence inside $D$, while extensivity gives the reverse
inclusion. Their complements are consequently a topology. Now $c(A)$ is
fixed and is contained in every fixed supermask of $A$, so it is exactly
the topological closure. The specialization matrix can be read directly
from $x\in c(\{y\})$.

\begin{proposition}[A minimum dense subset]\label{idx-finite114}
Identify points $x,y$ when $x\preceq y\preceq x$. A mask is dense exactly
when it meets each maximal equivalence class in the resulting finite
partially ordered set. Choosing the least numbered point in each maximal
class gives a dense subset of minimum cardinality.
\end{proposition}
\begin{proof}
The relation on classes is well defined by transitivity. Above any class
one can find a maximal class by moving to the least labelled strict
successor as long as one exists. A strict ascent cannot revisit a class
by antisymmetry. More explicitly, the number of classes strictly above
the current one drops at each move. Iterating $n$ times, with terminal
points fixed, suffices. If $A$ meets all maximal classes, the chosen
point of a maximal class above $x$ witnesses $x\in\operatorname{cl}A$.
Conversely a witness above a point of a maximal class must belong to that
class. Since the classes are disjoint, their chosen representatives give
the lower bound as well as a set attaining it (an injection into every
dense set selects its first member in each such class).
\end{proof}

\subsection{Maps, separation and classification by open sets}

\begin{theorem}[Continuity tests]\label{idx-finite136}
For a table $f:X\to Y$, each of the following characterizes continuity:
monotonicity for specialization, preservation of closed sets under inverse
image, and $f(\operatorname{cl}A)\subseteq\operatorname{cl}(f(A))$ for
every mask $A$. A bijection is a homeomorphism exactly when it preserves
and reflects the relation. Homeomorphism of two given finite spaces is
primitive-recursively decidable.
\end{theorem}
\begin{proof}
A monotone map takes upper sets to upper sets under inverse image. For the
converse, $f^{-1}(N_{f(x)})$ is an open through $x$, so it contains every
$y$ with $x\preceq y$. Thus $f(x)\preceq f(y)$. The closed-set test follows
by complementing preimages. Given monotonicity, carry a closure witness
$x\preceq y\in A$ to $f(x)\preceq f(y)\in f(A)$. Conversely apply the
closure condition to $A=\{y\}$ for every $x\preceq y$. Identities and
composites preserve the relation. A bijection has monotone inverse precisely
when it reflects the relation; its inverse table uses first-preimage search.
One can enumerate all forward tables on the carriers and retain exactly
those with the bijection and two relation tests. For domain size $n$ and
target size $m$, digits with radix $m+1$ give bound $(m+1)^n$ after checking
each digit is $<m$. This convention handles the empty carriers as well.
\end{proof}

\begin{theorem}[Separation, reflection, and the Sierpi\'nski space]\label{idx-finite159}
A finite space is $T_0$ exactly when its preorder is antisymmetric. The
conditions $T_1$, Hausdorff, and discrete are equivalent here. The quotient
identifying mutually comparable-in-both-directions points is $T_0$ and
is universal for maps to $T_0$ spaces. For finite $X$, opens of $X$ are in
bijection with continuous maps to the two-point space
$S=(\{0,1\},0\preceq 1)$. Every finite $T_0$ space embeds in a finite
power of $S$.
\end{theorem}
\begin{proof}
The $T_0$ test asks for an open containing exactly one of two distinct
points. When $x\not\preceq y$, $N_x$ is one; if both comparisons hold every
open has equal membership at $x,y$. All comparison bits are decidable.
The $T_1$ test requires an open through $x$ missing $y$ for each ordered
pair of distinct points, so $N_x=\{x\}$ for all $x$ and every mask is open.
Discrete spaces separate distinct points by disjoint singletons. Disjoint
neighbourhoods imply the $T_1$ test.

Represent the equivalence class of $x$ by $r(x)$, the least $y$ with
$x\preceq y\preceq x$, and retain the subset $R$ of representatives. For
representatives the restricted relation is antisymmetric. For all points
$x\preceq y$ is equivalent to $r(x)\preceq r(y)$ by transitivity, so
projection $r$ and inclusion $j:R\to X$ are monotone. Any monotone $f$
to a $T_0$ target takes equivalent points to equal points by antisymmetry.
The unique factor is $f$ restricted to $R$ because $rj$ is the identity.

Finally $\chi_U:X\to S$ is monotone precisely when $U$ is upper. Use the
coordinates $e(x)_a=\chi_{N_a}(x)$ for $a\in X$. Coordinatewise order
on the tuples gives $e(x)\le e(y)$ exactly when $x\preceq y$: monotonicity
gives one direction, and the coordinate $a=x$ gives the other. Antisymmetry
makes $e$ injective. The induced order on the image is the subspace order
by the construction below. This proves the embedding, including the empty
space.
\end{proof}

A separate example marks the role of the inverse-continuity hypothesis.
The identity from the discrete two-point space to $S$ is a continuous
bijection; its inverse fails monotonicity at $(0,1)$. For finite spaces a
continuous bijection to a Hausdorff target does have continuous inverse:
the inverse images of its open singletons are singleton opens throughout
the domain.

\subsection{Reconstructing the points from the lattice of opens}

For upper masks $U,V$, meet and join are intersection and union. Their
Heyting implication also has a finite formula:
\[
 U\Rightarrow V=\{x:(\forall y\in N_x)\,[y\in U\ \Longrightarrow\ y\in V]\}.
\]
This is upper since $x\preceq z$ implies $N_z\subseteq N_x$. For any open
$W$, if $W\subseteq U\Rightarrow V$ then $W\cap U\subseteq V$ by using
$y=x$ at each $x\in W\cap U$. Conversely, if $W\cap U\subseteq V$ and
$x\in W$, every $y\in N_x$ lies in $W$ and so $y\in U$ forces $y\in V$.
This proves the Heyting adjunction. In the two-point example the
pseudocomplement $\{1\}\Rightarrow\varnothing$ is empty, so the join of
this open with its pseudocomplement is a proper open.

\begin{theorem}[Finite distributive lattice representation]\label{idx-finite216}
Given the operation tables of a finite bounded distributive lattice $L$,
iPRA constructs a finite $T_0$ space whose open-set lattice is isomorphic
to $L$. Conversely the open-set lattice determines a finite $T_0$ space
and all its continuous maps: every bounded lattice homomorphism
$h:\mathcal O(Y)\to\mathcal O(X)$ between open-set lattices of finite
$T_0$ spaces is the inverse-image operation of a unique continuous
$f:X\to Y$.
\end{theorem}
\begin{proof}
Set $a\le b$ when $a\wedge b=a$. Scan the table for the nonzero elements
$j$ such that $j=u\vee v$ always entails $j=u$ or $j=v$ (the
join-irreducibles). Every such $j$ is join-prime: if $j\le a\vee b$,
distributivity gives $j=(j\wedge a)\vee(j\wedge b)$, hence $j\le a$ or
$j\le b$. Induction on the length extends this test to a finite join,
including the empty join since $j\ne0$.

Each $a$ is the join of all join-irreducibles below it. For $0$ this is the
empty join. If $a$ is join-irreducible, it is itself among the factors.
Otherwise search the operation table: since all the tests are decidable,
for $a\ne0$ there are $u,v$ with $a=u\vee v$ and $u,v<a$.
Each of these has fewer elements below or equal to it than $a$ does, since
its lower mask omits $a$. Induct up to $|L|$ on this lower-mask cardinality,
with the invariant testing the asserted join equation for all elements of
cardinality at most the counter. The joins for $u,v$ then give $a$; adding
any further irreducibles below $a$ cannot increase it beyond $a$.

Let $P$ consist of the join-irreducibles, ordered oppositely to $L$.
Then $a\mapsto D_a=\{j:j\le a\}$ is an upper mask on $P$. It preserves
meet by the definition of lower bound, and join by join-primality. It
preserves both bounds. It reflects equality by the join formula. If $D$ is
an upper mask of $P$, put $a=\bigvee D$ in $L$. Then $D\subseteq D_a$;
conversely any $j\in D_a$ is below some member of $D$ by the finite
join-prime test, so lies in $D$. This gives a computable inverse, also
when $L$ is the one-element lattice and $P$ is empty.

For the map statement and $x\in X$, intersect all opens $U$ of $Y$ with
$x\in h(U)$ to obtain an open $B$. This family includes $Y$, and
$x\in h(B)$ since $h$ preserves finite meets. Because $h$ preserves the
empty open, $B$ is inhabited. Write $B$ as the union of the $N_y$ over
$y\in B$. Preservation of finite joins gives a first $y\in B$ with
$x\in h(N_y)$. By the defining intersection $B\subseteq N_y$, and
$N_y\subseteq B$ since $B$ is open. Thus $B=N_y$ and antisymmetry
uniquely determines $y$ from $B$. Put $f(x)=y$. If $x\in h(U)$, then
$y\in B\subseteq U$. If $y\in U$, then $B=N_y\subseteq U$, so
$x\in h(B)\subseteq h(U)$; monotonicity of $h$ follows from meet
preservation. Thus $h(U)=f^{-1}(U)$ on every open, giving continuity.
For a second such function, the membership of its value in every open
would be identical, and $T_0$ forces the values equal. Applied to identity
and composite homomorphisms, uniqueness proves the stated reconstruction
up to isomorphism: in particular two isomorphic lattices give mutually
inverse continuous maps. All joins, meets, tests for irreducibles and
first-witness searches here scan bounded tables.
\end{proof}

\subsection{Building spaces and testing their universal properties}

\begin{theorem}[Subspaces, products, sums and quotients]\label{idx-finite273}
The induced topology on a mask $A\subseteq X$ corresponds to the restricted
preorder. The product preorder is coordinatewise, and the disjoint sum
preorder compares only points of the same summand. For a surjection
$q:X\twoheadrightarrow Q$, the quotient topology is given by the reflexive
transitive closure of
\[
 uEv\quad\Longleftrightarrow\quad
 (\exists x,y\in X)\,[q(x)=u,\ q(y)=v,\ x\preceq y].
\]
All four constructions and their map tests are primitive recursive.
\end{theorem}
\begin{proof}
For a point of $A$, its least relative neighbourhood is $A\cap N_x$.
Indeed any relative open through $x$ contains this intersection, which
is itself relative open. For the product, the rectangles of principal
neighbourhoods are precisely the principal upper sets of the coordinatewise
preorder. Every upper mask is a union of such rectangles, and every product
of opens is upper; this identifies the topology generated by open rectangles.
A table into the product is monotone exactly when each of its two projections
is monotone. In the sum the upper masks restrict to arbitrary upper masks
on the summands, proving the usual open-set test, and a table out of the
sum is monotone precisely when both restrictions are.

In the quotient, $q^{-1}(V)$ is upper if and only if membership of $V$
propagates along every $E$-edge. This holds if and only if it propagates
along every finite directed $E$-path by induction along that path. The
reflexive transitive closure can be computed by starting with the identity
and adding a last edge at each of $|Q|$ rounds, retaining all old pairs.
The invariant states that after $k$ rounds entries represent paths of
at most $k$ edges. Repeated vertices in a path allow the intervening
segment to be erased without changing its endpoints. Alternatively,
from each start the reachable mask grows strictly until stationary, so is
stationary by $|Q|$ rounds. A predecessor table at each round provides a
path for every positive entry. Therefore the ultimate matrix is transitive
and has exactly the upper masks prescribed by the quotient open-set test.
A function $h:Q\to Z$ is monotone exactly when $hq$ is: necessity uses
monotonicity of $q$, and sufficiency propagates monotonicity from the
generating edges along paths. In particular a monotone $f:X\to Z$ constant
on fibres factors uniquely; at $u$ use the first $x$ with $q(x)=u$.
\end{proof}

These rules also yield the pasting lemma for a coded finite cover of $X$
by open masks, or by closed masks. Given continuous functions on the
pieces which agree on intersections, take at $x$ the value on the first
piece containing it. Agreement makes this the common value on any such
piece. To prove monotonicity for $x\preceq y$, an open piece containing
$x$ contains $y$; in the closed-cover case a piece containing $y$ contains
$x$. On that piece use its monotonicity. This proves continuity and the
restriction property, which in turn gives uniqueness.

\begin{proposition}[Finite function spaces]\label{idx-finite324}
Let $Y^X$ be the mask of monotone tables $X\to Y$ among the radix-coded
candidates above, with $f\preceq g$ when $f(x)\preceq g(x)$ for each $x$.
Evaluation is continuous. For each finite $Z$ there is a bijection
between continuous maps $Z\times X\to Y$ and $Z\to Y^X$. This topology
on $Y^X$ agrees with its compact-open topology for finite $X$.
\end{proposition}
\begin{proof}
If $(f,x)\preceq(g,y)$, then $f(x)\preceq g(x)\preceq g(y)$; this
proves evaluation is monotone. For $a:Z\times X\to Y$ monotone, its row
$x\mapsto a(z,x)$ is monotone, and rows increase pointwise with $z$.
Conversely evaluating such a table of rows produces a monotone $a$.
Digit lookup shows the two operations undo each other. For the last
assertion a compact-open subbasic set consists of the functions sending
a specified compact mask $K$ into a specified open $U$. It is upper.
Conversely the principal upper set through $f$ is the intersection over
$x\in X$ of these subbasic sets for $K=\{x\}$ and $U=N_{f(x)}$.
Singletons are compact since any specified cover of one has a containing
member. Hence all principal upper sets belong to the compact-open topology.
\end{proof}

\subsection{Components, compactness and finite covers}
\label{sec:finite-components}

\begin{theorem}[Components and clopens]\label{idx-finite348}
Join two points by an undirected edge when either specialization comparison
holds. Reachability in this graph computes the connected components and
the quasicomponents. It produces either a zigzag between two input points
or a specific clopen mask separating them. Every clopen is a union of
components; if there are $c$ components there are exactly $2^c$ clopens.
\end{theorem}
\begin{proof}
Use the reachability algorithm of the quotient construction with the
symmetric edge relation. Its classes are both upper and lower: an
edge from or to a class member meets the same class. Each is therefore
clopen. Within a reachable class, any clopen for the induced topology is
both upper and lower along edges inside it. Its membership bit is constant
along each witnessing zigzag and hence across the class. This proves each
class is connected (has no nonempty proper relative clopen). Any connected
mask containing $x$ lies in its class, since intersection with the class
and its complement would otherwise give a nontrivial clopen partition.
The classes are therefore the maximal inhabited connected masks. The
quasicomponent at $x$ is the intersection of the clopens through $x$.
Such clopens contain the class, and the class itself is among them.

Assign disjoint indices to classes by their least representatives. The union
of the classes indexed by any bit mask is clopen and all clopens are
obtained exactly once, by testing their membership at each representative.
The empty space has no components and has its one clopen; it has no
separation into two inhabited opens. These conventions also show that a
space has no separation if and only if every pair of points is joined by
a zigzag.
\end{proof}

Continuous images of connected masks are connected: a monotone table sends
a zigzag in the restricted preorder to a zigzag inside the image. A product
of two connected inhabited spaces is connected: follow the first coordinate's
zigzag with second coordinate fixed, then follow the second coordinate's
zigzag with first coordinate fixed. If either carrier is empty the product
has no separation. In particular the clopen count and component count are
homeomorphism invariants. For a nonempty family of connected masks with a
given common point $p$, their union is connected, since each of its points
is joined inside its own mask to $p$. Here the family is a table and $p$
is part of the hypotheses.

\begin{proposition}[A sharp subcover bound]\label{idx-finite389}
For a finite space $X$, let $b$ be the number of minimal specialization
classes after passage to $T_0$. From any list of open masks covering $X$ one
can extract a subcover of length at most $b$; this bound is attained by
some cover. Every subspace satisfies the same theorem in its restricted
order. Closed subsets have the corresponding finite-intersection test.
\end{proposition}
\begin{proof}
Take the first representative of each minimal class and, by bounded search,
the first member of the list covering it. Each point lies above a minimal
class by the descending version of the strict-ascent argument. The selected
opens therefore cover; repetitions may be removed. The principal opens of
the minimal representatives form a cover, and none contains a different
minimal representative, so every member of that cover is indispensable.
If $X$ is empty the empty list suffices. Apply the same argument to the
relative preorder for any subspace. Complementing a list of closed masks
with empty intersection yields an open cover, and the selected at most
$b$ complements correspond to closed masks with empty intersection.
\end{proof}

For an indexed open family supplied by a fixed primitive-recursive
membership predicate together with a primitive-recursive selector assigning
a containing index to each point, the same proof takes the indices supplied
at the minimal representatives. For a finite table the selector is computed
from the table. This spells out which form of finite-space compactness
the arithmetic proof establishes, even when the outside indexing set is
presented as unbounded.

\subsection{A classification theorem for finite homotopy certificates}

A \emph{zigzag homotopy} from $f$ to $g$ is a list of monotone maps
$f=h_0,\ldots,h_s=g$ such that each consecutive pair is pointwise
comparable in one direction on the entire domain. This is a numerical
certificate: it equivalently gives a monotone map on $X$ times a finite
fence, the latter generated by an oriented edge between each pair of
successive list positions. The monotonicity check follows by changing
one coordinate at a time. Conversely restriction to the successive positions
gives the original certificate. Composition on either side preserves these
certificates. Reversing and concatenating lists proves symmetry and
transitivity.

There is a geometric interpretation of a certified comparison $f\preceq g$
in ordinary topology. Use the homotopy taking value $f$ at time $0$ and
$g$ at positive times. For an open $U$ its inverse image is
\[
  (f^{-1}U\times[0,1])\ \cup\ (g^{-1}U\times(0,1]),
\]
which is relatively open, and the restriction to either endpoint is as
claimed. This describes the inverse images directly, with no decision about
equality of a real name required by the certificate verifier. Reverse or
concatenate these homotopies for a whole zigzag. The arithmetic
classification below concerns the specified finite certificates. To see the
relation in the other direction in the usual interpretation, let $H$ be an
ordinary continuous interval homotopy between maps of finite spaces. Its
slices define a continuous map $s:[0,1]\to Y^X$: at any fixed $t_0$ the
inverse image of the least neighbourhood of $s(t_0)$ is the finite
intersection, over $x\in X$, of the open time sets on which
$H(x,t)\in N_{H(x,t_0)}$. Since each component of $Y^X$ is clopen,
the inverse image of the component of $s(0)$ is clopen and contains $0$.
Connectedness of the ordinary interval puts $s(1)$ in that component too.
The finite reachability algorithm then returns a zigzag. This comparison
uses the ordinary interval and its connectedness externally; the arithmetic
core proof below has only numerical inputs.

\begin{theorem}[Core reduction and rigidity]\label{idx-finite453}
In a finite poset call $x$ an upper beat point if its strict upper set has a
least element, and a lower beat point if its strict lower set has a greatest
element. Compute a core by passing to the $T_0$ representatives and deleting
beat points until none remain. Any such core is a retract of the given space
whose retraction followed by inclusion is zigzag-homotopic to the
identity. Two finite spaces admit mutually inverse maps up to zigzag homotopy
if and only if their cores are order isomorphic. This gives a
primitive-recursive decision and produces equivalence witnesses on success.
\end{theorem}
\begin{proof}
For the first reduction use $r,j$ from the separation theorem.
We have $rj=1$ and $jr$ is pointwise comparable to $1$ in both directions.
If $x$ is an upper beat point with least strict successor $y$, let $r$
fix the other points and send $x$ to $y$. Check an ordered pair
$a\preceq b$. If $a=x\ne b$, minimality gives $y\preceq b$; if
$b=x\ne a$, use $a\preceq x\preceq y$. Equality and the cases avoiding
$x$ follow at once. Thus $r$ is monotone, restricts to the identity on
the remaining subposet, and satisfies $1\preceq jr$. The lower case is
obtained by reversing every order comparison. The single comparison fixes
the remaining points. For successive deletions, composing the homotopy from
each step with the earlier inclusions and retractions gives the required
list; each earlier inclusion acts as identity on the ultimate subspace.
Tests for beat points and their witnesses scan pairs and all intermediate
points. Delete the least eligible point, preferring an upper deletion if
both tests pass, and iterate at most $n$ times. Alternatively any supplied
sequence of eligible choices gives the same properties.

Let $P$ be a beat-free finite poset and $h:P\to P$ monotone with
$1\preceq h$. If $h$ moves a point, choose a maximal moved point $x$ by
bounded ascent. For every $z>x$, maximality gives $h(z)=z$, so
$x<h(x)\preceq h(z)=z$. Then $h(x)$ is the least strict successor of
$x$, a contradiction. Because the set of moved points is a mask this proves
$h=1$. Dually $h\preceq 1$ also forces $h=1$. Along a list starting at
$1$, induction gives that every subsequent map is $1$. The prefix
invariant tests equality of two finite tables.

Given homotopy inverses between two original spaces, insert the
inclusions and retractions on both sides to get maps between their cores.
Their composites are zigzag-homotopic to the identities by composing the
certificates just constructed with the supplied maps. Rigidity makes both
composites actual identities. Thus these core maps are order isomorphisms.
Conversely, compose any isomorphism between two cores and its inverse with
the original retractions and inclusions. The displayed retraction
certificates verify the homotopy-inverse equations. The previous bounded
isomorphism search, followed by these explicit compositions, proves the
decision statement. Taking two reduction sequences on the same input gives
uniqueness of their cores up to isomorphism.
\end{proof}

For example, on a nonempty finite poset with least element the identity is
pointwise above the constant map to that element, so the space contracts in
one comparison. In contrast, take two incomparable minimal points each below
each of two incomparable maximal points. This four-point poset is connected
by zigzags. It has no beat points, since every nonempty strict upper or
lower set has two incomparable members. Its core has four points and
hence cannot be isomorphic to a singleton. It is connected and fails to
be contractible in the certificate model.

\begin{corollary}[Fixed points of finite-space maps]\label{idx-finite512}
A finite $T_0$ space has the property that every continuous self-map has a
fixed point if and only if its beat-free core has that property. In
particular every continuous self-map of a finite $T_0$ space that is
zigzag-contractible to a point has a computable fixed point.
\end{corollary}
\begin{proof}
For one deletion write $i:Y\to X$ for inclusion and $r:X\to Y$ for the
retraction. If $X$ has the property, apply it to $igr$ for $g:Y\to Y$.
Any fixed point belongs to $Y$ and there is fixed by $g$. Conversely,
suppose $Y$ has the property and let $f:X\to X$. Find a fixed point $y$
of $rfi$ by scanning $Y$. If $f(y)=y$ we are done. Otherwise $f(y)$
must be the deleted point $x$, and $y=r(x)$. For an upper deletion
$x<y$, so $f(x)\le f(y)=x$. Iterating $f$ from $x$ gives a descending
chain. Whenever a step changes value, the number of elements below or
equal to it strictly decreases, so at most $|X|$ iterations give a
stationary point. For a lower deletion reverse the inequalities.
This proves equivalence through the entire reduction. A contractible
$T_0$ space has singleton core by the rigidity theorem; the unique
self-map there has a fixed point. The $T_0$ assumption matters: both
points of the indiscrete doubleton are topologically indistinguishable,
and their transposition is a continuous fixed-point-free map although
the space contracts to either point by a single comparison.
\end{proof}

\subsection{Lean construction for the finite theorems}

At the numerical layer use \texttt{Nat} for each bounded table and use
\texttt{bit(n,A,i)} for membership, taken to be false when $i\ge n$.
Use \texttt{allBelow} and \texttt{anyBelow} by recursion on their bound.
Tabulate a predicate $p$ on $[k]$ by summing $2^i$ for the true indices.
An $n$-row relation uses bit $nx+y$, and a map is read by digit extraction
in radix $m+1$, with an explicit digit-range test. The following dependency
order gives the open proof obligations; names describe suggested Lean
declarations.
\begin{enumerate}
\item Prove \texttt{bit\_tabulate}, prefix specifications for the two
bounded scans, first-witness selection (with the bound as failure sentinel),
and digit-table extensionality for normalized codes. Define
\texttt{IsTopology}, \texttt{IsPreorder}, \texttt{upset}, \texttt{principal},
\texttt{specializes} and \texttt{opensOf} by the tests above. The invariant
for the defining intersection says a bit survives iff it survived every
selected mask in the processed prefix. This yields both round trips
by bit extensionality.
\item Define closure by witness search and interior with \texttt{allBelow}.
Prove the local formulas first. Then closure idempotence is a single use
of relation transitivity between the two chosen witnesses. Define
\texttt{MonotoneTable} by scanning ordered pairs. Use the mask
\texttt{principal (f x)} for the substantive direction of
\texttt{continuous\_iff\_monotone}.
\item Implement \texttt{reachStep} by adding endpoints of single edges
from the previous mask; retain for new endpoints their first predecessor.
Induct on the round to show that all recorded predecessors recover paths.
The cardinality of the current mask strictly increases if a new bit
appears; prove this by a scan counting bits. If a round adds nothing,
subsequent rounds repeat it. This proves stationarity at the carrier
bound. Applying the same routine to symmetric and directed edge tables
gives respectively the component and quotient constructions.
\item Use the least index in each equivalence class for the reflection,
count representatives, and implement product indexing by mixed radix.
Define currying and uncurrying by direct table lookup and prove their
equations entrywise. Select each cover index with first-witness search;
the minimum and maximum class searches use at most $n$ order steps.
\item For the lattice theorem, form the table of join-irreducibles by testing
all pairs. At stage $k$ verify the join reconstruction equation for all
$a$ whose lower mask has cardinality at most $k$. For the reconstruction
of a lattice homomorphism use the bounded intersection of qualifying opens
and extract the least principal open equal to this intersection.
\item Represent a homotopy certificate by a coded list of normalized
maps of the same dimensions. Verify each map and each adjacent comparison,
including equality at the endpoints. Implement \texttt{coreStep} on a mask
of surviving points, testing strict successors and predecessors only there.
Its invariant keeps a retraction, the inclusion and a list of maps from
the original identity to their composite. One deletion appends the new
composite, which is comparable to the preceding composite. The proof of
rigidity needs the bounded maximization of a moved point and induction
along the supplied list. Isomorphism of final cores is a search over maps,
with inverses computed on the surviving points or after increasing-order
reindexing. None of these loops needs induction on a formula with an
unbounded quantifier.
\end{enumerate}
